\documentclass[11pt]{article}
\usepackage[margin=1in]{geometry}
\usepackage{amsmath,amssymb,amsthm}
\usepackage[hidelinks]{hyperref}
\usepackage{microtype}

\newtheorem{theorem}{Theorem}
\newtheorem{lemma}{Lemma}
\newcommand{\R}{\mathbb R}
\newcommand{\PhiC}{\Phi}

\title{A Literature-Implied Answer to the $\Phi_1\setminus\Phi_2$
Negative-Squares Conjecture}
\author{Run \texttt{fa\_banach\_001}: candidate status note}
\date{13 August 2026}

\begin{document}
\maketitle

\noindent\textbf{Status.}
\emph{Literature-implied answer (full conjecture; candidate likely valid).}
The supporting theorem predates the conjecture and does not mention it.  The
short implication below appears not to have been identified in the bounded
search described at the end of this note.

\section{Source question}

Golinskii, Malamud, and Oridoroga define, for a continuous real function
$f$ on $[0,\infty)$,
\[
 \kappa_n^-(f)=\sup_X n_-\bigl([f(|x_i-x_j|)]_{i,j}\bigr),
\]
where $X$ ranges over finite subsets of $\R^n$ and $n_-$ counts negative
eigenvalues with multiplicity.  On PDF page 3 of
\cite{GMO} they state:
\begin{quote}
For each function $f\in\Phi_1\setminus\Phi_2$, the relation
$\kappa_2^-(f)=+\infty$ holds.
\end{quote}
Here $f\in\Phi_n$ means that $x\mapsto f(|x|)$ is positive definite on
$\R^n$.

\section{Supporting classification}

We use Theorem~3.2 of Sasvari \cite{Sasvari1985}, on journal pages
229--231 (PDF pages 7--9 of the local primary scan).  In the
special form needed here it says:

\begin{theorem}[Sasvari's bounded classification]\label{thm:sasvari}
Let $G$ be a locally compact abelian group.  If a bounded continuous
Hermitian function $h$ on $G$ has exactly $k<\infty$ negative squares,
then
\[
 h(x)=\int_{\widehat G}\gamma(x)\,d\mu_+(\gamma)
      -\int_{\widehat G}\gamma(x)\,d\mu_-(\gamma),
\]
where $\mu_+$ and $\mu_-$ are finite positive measures, $\mu_-$ is
concentrated at exactly $k$ distinct characters, and $\mu_+$ gives those
characters no mass.  Conversely, such a representation has exactly $k$
negative squares.
\end{theorem}

The introduction of Sasvari's later paper \cite[p.~319]{Sasvari1989}
independently restates this result: bounded functions with $k$ negative
squares are precisely Fourier transforms of measures that are negative at
$k$ points of the character group and nonnegative outside those points.

\section{The implication}

\begin{lemma}\label{lem:line}
Let $\sigma$ be a finite positive rotation-invariant Borel measure on
$\R^2$ with $\sigma(\{0\})=0$.  Then every line $L$ through the origin has
$\sigma(L)=0$.
\end{lemma}

\begin{proof}
If $\sigma(L)>0$, then $\sigma(L\setminus\{0\})>0$.  Choose countably many
rotations taking $L$ to pairwise distinct unoriented lines $L_j$.  The sets
$L_j\setminus\{0\}$ are pairwise disjoint and, by rotation invariance, all
have the same positive measure.  Their union would have infinite measure,
contrary to finiteness of $\sigma$.
\end{proof}

\begin{theorem}\label{thm:answer}
For every $f\in\Phi_1\setminus\Phi_2$ one has
$\kappa_2^-(f)=+\infty$.
\end{theorem}

\begin{proof}
Set
\[
 g(x)=f(|x|),\qquad x\in\R^2.
\]
Because $f\in\Phi_1$, the $2\times2$ positive-definiteness inequality on
the real line implies $|f(t)|\leq f(0)$; hence $g$ is bounded.  It is also
continuous and Hermitian.

Suppose, for a contradiction, that $\kappa_2^-(f)=k<\infty$.  Since
$f\notin\Phi_2$, we have $k\geq1$.  The supremum defining $k$ is over
integers, so some finite Schoenberg matrix has exactly $k$ negative
eigenvalues; thus $g$ has exactly $k$ negative squares in Sasvari's sense.
Identify $\widehat{\R^2}$ with $\R^2$.  Theorem~\ref{thm:sasvari} gives a
finite signed measure
\[
 \mu=\mu_+-\mu_-,\qquad
 g(x)=\int_{\R^2}e^{i\langle x,\xi\rangle}\,d\mu(\xi),              \tag{1}
\]
whose negative part $\mu_-$ is supported at exactly $k$ points.

For each planar rotation $R$, radiality gives $g(Rx)=g(x)$.  Uniqueness of
Fourier--Stieltjes transforms of finite measures in (1) implies
$R_*\mu=\mu$ (replacing $R$ by its transpose, which runs through the same
rotation group).  Uniqueness of the Jordan decomposition then gives
$R_*\mu_+=\mu_+$ and $R_*\mu_-=\mu_-$.  The support of $\mu_-$ is therefore
a finite rotation-invariant subset of $\R^2$.  Every nonzero point has an
infinite rotation orbit, so
\[
 \mu_-=a\delta_0\quad\text{for some }a>0.                           \tag{2}
\]
Moreover $\mu_+(\{0\})=0$, because $\mu_+$ and $\mu_-$ are the mutually
singular Jordan parts.

Let $\pi(\xi_1,\xi_2)=\xi_1$.  Restricting (1) to the first coordinate
axis shows that
\[
 f(|t|)=g(t,0)=\int_{\R}e^{its}\,d(\pi_*\mu)(s).
\]
The left side is positive definite on $\R$, since $f\in\Phi_1$.
Bochner's theorem and uniqueness of Fourier--Stieltjes transforms therefore
force the signed measure $\pi_*\mu$ to be positive.

On the other hand, $\pi^{-1}(\{0\})=L^\perp:=\{(0,s):s\in\R\}$.  Apply
Lemma~\ref{lem:line} to the finite rotation-invariant measure $\mu_+$:
$\mu_+(L^\perp)=0$.  Using (2),
\[
 (\pi_*\mu)(\{0\})
 =\mu(L^\perp)
 =\mu_+(L^\perp)-a
 =-a<0,
\]
contradicting positivity of $\pi_*\mu$.  Hence no finite $k$ is possible,
and $\kappa_2^-(f)=+\infty$.
\end{proof}

\section{Identification and scope}

The proof uses no regularity, decay, or support hypothesis beyond
$f\in\Phi_1\setminus\Phi_2$.  It therefore establishes the full conjecture,
not merely the additional-measure-hypothesis cases proved in the source
paper.  The mathematical input is nevertheless Sasvari's 1985 theorem; the
new content of this packet is the direct radial-symmetry and line-projection
identification.  For that reason its provenance is
\emph{literature-implied answer}, not an original full result.

A bounded novelty/status search was run on 13 August 2026 in the run's
registry, solution, attempt, and proof-gap indexes and on the web using the
exact source title, the notation $\kappa_2^-(f)$, variants of
``$\Phi_1$, $\Phi_2$, negative squares'', and Sasvari's title and DOI.  No
later paper explicitly claiming this conjecture as resolved was found.  That
negative search is bounded and is not a claim of exhaustive bibliographic
novelty.

\begin{thebibliography}{9}
\bibitem{GMO}
L. Golinskii, M. Malamud, and L. Oridoroga,
\emph{Radial positive definite functions and Schoenberg matrices with
negative eigenvalues}, arXiv:1502.07179 (2015), conjecture on PDF p.~3.

\bibitem{Sasvari1985}
Z. Sasvari,
\emph{On bounded functions with a finite number of negative squares},
Monatsh. Math. \textbf{99} (1985), 223--234,
doi:\href{https://doi.org/10.1007/BF01295156}{10.1007/BF01295156};
Theorem~3.2, pp.~229--231.

\bibitem{Sasvari1989}
Z. Sasvari,
\emph{Characterization of locally bounded functions with a finite number of
negative squares}, Acta Sci. Math. (Szeged) \textbf{53} (1989), 319--327;
classification restated on p.~319.
\end{thebibliography}

\end{document}
